\documentclass[11pt,a4paper]{article}
\usepackage[T1]{fontenc}
\usepackage[utf8]{inputenc}
\usepackage{lmodern}
\usepackage[margin=26mm]{geometry}
\usepackage{amsmath,amssymb,amsthm,mathtools}
\usepackage{booktabs,array}
\usepackage{needspace}
\usepackage[hidelinks]{hyperref}
\newtheorem{theorem}{Theorem}[section]
\newtheorem{lemma}[theorem]{Lemma}
\newtheorem{corollary}[theorem]{Corollary}
\theoremstyle{definition}
\newtheorem{definition}[theorem]{Definition}
\theoremstyle{remark}
\newtheorem{remark}[theorem]{Remark}
\newcommand{\NN}{\mathbb N}
\newcommand{\QQ}{\mathbb Q}
\newcommand{\bits}{\operatorname{bits}}
\newcommand{\Time}{\operatorname{Time}}
\newcommand{\pa}{\operatorname{par}}
\newcommand{\RC}{\operatorname{rc}}
\title{Tree-budget search: parameter corollaries and a paper proof of bit-cost bounds}
\author{UAV routing certification project}
\date{5 November 2026}
\begin{document}
\maketitle
\begin{abstract}
We specify a binary-encoded matrix decision and explicit-witness task, and
give a paper construction of an explicit-vector implementation of the existing
ordered, association-list memoized search. Including candidate enumeration,
legality tests, both lookup and insertion key comparisons, arithmetic, and
witness construction, its bit-cost RAM time is
\(O(S^3(M+1)(\ell+1)^5)\).
Here \(S=\prod_{v\ne o}(H_v+1)\), \(M=\sum_{v\ne o}H_v\), and \(\ell\)
is the total input encoding length. This gives paper-level FPT upper bounds
for \(M\) and for the pair \((t,U)\). The representation simulation and the
bit-cost theorem in this note are not Lean theorems and are not a time bound
for the current compiled Lean program. Existing Lean results certify the
search semantics and logical operation counts used as premises.
\end{abstract}

\section{Status and scope}
The source baseline is the research commit
\href{https://github.com/huam54925-cpu/uav-routing-certification/commit/1f4774cfcb59c554893d37a9e4d476b0d2c8fa78}{\texttt{1f4774c}}.
Its modules \texttt{TreeBudget\allowbreak SearchStatsBounds},
\texttt{TreeBudget\allowbreak CandidateCount}, \texttt{TreeBudget\allowbreak SearchCostBounds}, and
\texttt{TreeBudget\allowbreak StateParameters} supply the logical bounds below.
The accompanying \texttt{TreeBudget\allowbreak ParameterizedBounds} module closes the
logical parameter corollaries; its exact build and audit status is recorded
separately in the repository status document.

The subsequent representation package verifies array-backed vector operations,
signed integer scaling (including unreduced raw fractions), Boolean row-check
equivalence, and bounds on every successful entry of a valid, duplicate-free,
reachable cache. These are component correspondences. Ordered candidate-list
correspondence, full search simulation, and a bit-cost semantics remain outside
the Lean results in this package.

Sections~\ref{sec:task}--\ref{sec:time} are mathematical proofs about an
explicit representation specified here, not claims of a completed Lean
refinement or compiler-cost proof. No performance measurements are used.
In particular, an arbitrary value of type \texttt{Fin n -> Nat} is not
assumed to have constant-time evaluation. The task starts from multiplicities
\(H\); converting a different demand encoding into \(H\), expanding rows
into physical tours, and compressing certificates are outside this task.

\section{Logical parameter corollaries}
Let \(V=\{0,\ldots,n-1\}\), where \(n\ge1\), be the vertices of a rooted
tree with root \(o\) and parent map \(\pa\). Set
\[
 S=\prod_{v\ne o}(H_v+1),\quad M=\sum_{v\ne o}H_v,\quad
 t=|\{v\ne o:H_v>0\}|,\quad
 U=\max(\{H_v:v\ne o\}\cup\{0\}).
\]
Write \(k_0=\RC(H)=\sum_{v\ne o,\,\pa(v)=o}H_v\).
For the formal cold entry, let \(P,E,Q,C,I,L\) denote the candidate-list
length, expansions, candidate checks, state calls, insertion scans, and lookup
key comparisons, respectively. The established bounds are
\begin{equation}\label{eq:logical}
 P,E\le S,\quad Q\le S^2,\quad C\le1+S^2,\quad
 I\le\frac{S(S-1)}2,\quad L\le S+S^3.
\end{equation}
The final cache has exactly \(E\) entries. These initial-entry bounds do not
apply to an arbitrary hot cache or arbitrary candidate list.

\begin{lemma}\label{lem:parameters}
We have \(k_0\le M\le tU\), \(S\le2^M\), and \(S\le(U+1)^t\).
A successful cold-entry row witness has exactly \(k_0\) rows, hence at most
\(M\) rows.
\end{lemma}
\begin{proof}
The root children form a subset of the nonroot vertices, and all summands
are nonnegative, giving \(k_0\le M\). Delete zero summands from \(M\) and
bound each of the remaining \(t\) summands by \(U\).
Induction gives \(h+1\le2^h\) for every \(h\in\NN\): the successor step
uses \(h+2\le2(h+1)\). Multiplication over nonroot vertices gives
\(S\le2^{\sum H_v}\). Removing unit factors from \(S\) leaves \(t\)
factors, each at most \(U+1\). Finally, the residual-witness specification
at remaining count \(k\) requires a list of length exactly \(k\).
Apply soundness at \(k=k_0\). Empty products equal one, so the argument
also covers \(M=t=U=k_0=0\).
\end{proof}

\begin{corollary}\label{cor:envelope}
For either \(A=2^M\) or \(A=(U+1)^t\),
\[
 P,E\le A,\quad Q\le A^2,\quad C\le1+A^2,\quad
 I\le\frac{A(A-1)}2,\quad L\le A+A^3.
\]
\end{corollary}
\begin{proof}
Substitute \(S\le A\) in \eqref{eq:logical}. For the insertion bound,
use monotonicity of \(\sum_{j=0}^{a-1}j=a(a-1)/2\) on natural numbers.
This avoids treating truncated subtraction as an independently monotone
expression. These are logical-count conclusions, without a time claim.
\end{proof}

\section{A fixed encoded task}\label{sec:task}
\begin{definition}[Input and output]\label{def:task}
The input consists of an explicitly listed parent array on \(V\), a root
index \(o\), a list of nonnegative binary integers \(H_v\), signed rational
service fees \(s_v=a_v/b_v\), and a rational budget \(T=a_T/b_T\), where
all denominators are positive integers. Fractions need not be reduced.
Let \(\ell\ge n\) be the total length of this encoding, including delimiters
and signs. The parent array is required to describe a rooted tree:
\(\pa(o)=o\), and every other vertex reaches \(o\) by following parents.
Malformed encodings may be rejected before search.

Replace \(H_o\) by zero after reading the input, and call the normalized
vector \(H\). The task is to decide whether there exists a list of
\(k_0=\RC(H)\) rows \(p\in\NN^n\), each satisfying
\begin{align}
 &p_o=0,\qquad 0\le p_v\le H_v\quad(v\ne o),\qquad\RC(p)=1,
       \label{eq:rowa}\\
 &v\ne o,\ p_v>0,\ \pa(v)\ne o\ \Longrightarrow\ p_{\pa(v)}>0,
       \label{eq:rowb}\\
 &\sum_{v\ne o}\bigl(s_v+\mathbf1_{\pa(v)\ne o}s_{\pa(v)}\bigr)p_v\le T,
       \label{eq:rowc}
\end{align}
whose nonroot column sums are \(H_v\). On success the witness task outputs
that explicit row matrix in binary; on failure it outputs \textsc{no}.
The decision task may run the same algorithm and discard the witness.
\end{definition}

These conditions are exactly \texttt{RowValid} and the initial matrix
condition, using \texttt{cost\_eq} for the linear service cost. Condition
\eqref{eq:rowb} is a local parent test, not an enumeration of ancestor paths
or subsets. No sign restriction on \(s_v\) or \(T\) is needed.
An optional binary row-count input \(m\) can be handled by rejecting
\(m\ne k_0\): summing \(\RC(p)=1\) over rows and using the column sums
forces \(m=k_0\). This note does not analyze the separate demand-to-flow
wrapper or root-branch decomposition.

\begin{remark}[Why normalization matters]
The low-level logical count theorems allow an arbitrary root coordinate.
Matrix feasibility, however, ignores that coordinate whereas the residual
base test checks \(r=0\) at every coordinate. The mathematical entry here
uses \(H_o=0\), exactly as \texttt{searchMemo} normalizes its input. Without
this step, a nonzero root coordinate could incorrectly force failure for
the intended matrix task. The values \(S,M,t,U,k_0\) are unchanged by
normalization.
\end{remark}

\section{Representation and semantic correspondence}\label{sec:representation}
Use materialized length-\(n\) binary vectors, not unevaluated function
closures. Enumerate all vectors with \(p_o=0\) and \(0\le p_v\le H_v\)
in the coordinate order of \texttt{boundedVectors}, then stably filter by
\eqref{eq:rowa}--\eqref{eq:rowc}. Store the resulting ordered list once,
and reuse that same list in every recursive call.

The cache is an ordered association list of keys \((k,r)\) and answers.
An answer is either failure or a successful row list; a cache miss is a
third, distinct outcome. In particular, a cached failure is a hit and a
successful empty list is not a failure. Use these operations:
\begin{enumerate}
\item Search the cache from its head, comparing full keys, and return the
      first matching answer if present.
\item On a miss with \(k=0\), return success with the empty list exactly
      when every coordinate of \(r\) is zero.
\item On a miss with \(k>0\), scan the full candidate list in order.
      For each \(p\le r\), materialize \(r-p\) and call the solver with
      \(k-1\). A failing child passes its updated cache to the next
      candidate. A successful child returns \(p\) prepended to its rows
      immediately, without visiting the unused suffix.
\item After the miss has been evaluated, filter the entire returned cache
      to remove matching keys, then prepend the new key and answer.
      The filter is still executed when the parent key is known to be absent.
\end{enumerate}
Start with \((k,r)=(k_0,H)\) and the empty cache. Immutable vectors and
row lists may be shared by references, or copied; the conservative bound
below allows copying an entire answer whenever an answer is transferred.

\begin{lemma}[Paper simulation]\label{lem:simulation}
Interpret a stored vector \(x\) as the function \(v\mapsto x_v\).
Under this interpretation, the algorithm just specified returns the same
answer and ordered cache as \texttt{memoRows}, and its logical events
\(C,E,Q,L,I\) coincide with \texttt{memoRowsStat}.
\end{lemma}
\begin{proof}
First induct on the number of coordinates. The nested coordinate loops
produce precisely the list obtained by the range, flat-map and cons steps
of \texttt{boundedVectors}, in the same order. Array tests for
\eqref{eq:rowa}--\eqref{eq:rowc} express exactly \texttt{RowValid}; stable
filtering therefore gives the identical candidate list.

Next, array equality is extensional function equality on the finite domain.
Coordinatewise guarded subtraction agrees with natural-number subtraction,
and the all-zero array test agrees with function equality to zero.
Induct on \(k\), with a subsidiary induction on the scanned suffix.
Lookup visits identical entries and gives the same first hit or miss.
At zero depth both programs make the same zero test. At positive depth,
a rejected guard moves to the same suffix; an accepted guard invokes the
induction hypothesis on the same child state and input cache. On failure
the suffix induction starts with the child's returned cache. On success
both prepend the same row and stop. Finally the full-key filter and
prepending operation agree entry by entry.

Counting a state call, a miss, a candidate test, each visited lookup entry,
and each visited insertion-filter entry in these matching branches gives
the stated equality of events. This proves a mathematical simulation;
the complete explicit-vector search and this refinement have not been
encoded or checked in Lean in this package.
\end{proof}

\section{Bits, arithmetic, and the complete cost ledger}\label{sec:ledger}
Use a bit-cost RAM with binary storage and exact integer arithmetic.
Reading or writing a \(w\)-bit value or address is charged \(O(w)\),
not unit cost for unbounded words. Schoolbook addition, subtraction,
comparison, multiplication, and long division give \(O(w^2)\) upper bounds
for operands of at most \(w\) bits. Arrays are addressed by binary indices;
their address costs are included. Let
\[
 B=\ell+\lceil\log_2(M+2)\rceil+2.
\]
All constants in the following bounds are absolute and independent of
the parameters. We deliberately use loose polynomial bounds.

\begin{lemma}[Value and rational size]\label{lem:bits}
Every residual coordinate is between zero and \(H_v\), and every reachable
remaining count is at most \(k_0\le M\). All arithmetic needed for candidate
legality can be performed with \(O(B)\)-bit integers after polynomial-time
preprocessing.
\end{lemma}
\begin{proof}
The residual assertion follows by induction on guarded subtraction;
\(k\) only decreases. Nonroot coordinates are at most \(M\).
To avoid any repeated-denominator-growth assumption, set
\[
 D=b_T\prod_{v\in V}b_v>0,\quad
 A_v=a_v(D/b_v),\quad A_T=a_T(D/b_T),\quad
 c_v=A_v+\mathbf1_{\pa(v)\ne o}A_{\pa(v)}.
\]
The quotients are exact integers. The sum of the bit lengths of all input
denominators is at most \(\ell\), so \(D\) has \(O(\ell)\) bits. The
integers \(A_v,A_T,c_v\) have \(O(\ell)\) bits as well. Multiplication
by the positive \(D\) changes \eqref{eq:rowc} into exactly
\(\sum_{v\ne o}c_vp_v\le A_T\). For every partial sum, its absolute
value is at most \(\max_v|c_v|\,M\), hence its length is \(O(B)\).
This also handles negative fees, cancellations, and \(M=0\).
Forming the products and exact quotients uses \(O((n+1)B^2)\) bit work
under schoolbook algorithms. There is no rational arithmetic in the search
after candidates have been filtered.
\end{proof}

\begin{lemma}[Representation and address size]\label{lem:space}
An implementation as above can use addresses of \(O(B)\) bits throughout,
even if discarded allocations are never reclaimed. A key requires
\(O((n+1)B)\) bits, and any successful answer requires
\(O((M+1)(n+1)B)\) bits when copied explicitly.
\end{lemma}
\begin{proof}
The key claim follows from Lemma~\ref{lem:bits}. The cold execution preserves
cache validity, absence of duplicate keys, and reachability of every key.
Absence of duplicates makes each stored entry the answer visible to lookup;
lookup validity alone would not constrain answers hidden by duplicate keys.
Combining these three invariants gives \(|\mathrm{rows}|=k\le k_0\le M\)
for every stored successful answer, not merely the final result; each row
has \(n\) coordinates of \(O(B)\) bits. Candidate storage has \(P\le S\)
rows. Along the depth-first evaluation there are at most \(M+1\) state
frames; scans and cache traversals may instead be implemented by loops.

A finite implementation uses at most a fixed polynomial in
\(S,M+1,n+1,B\) elementary allocations: enumerate at most \(O((n+1)S)\)
partial vectors, perform at most \(C+Q+L+I+E+1=O(S^3)\) search events,
and permit a copy of at most \(M\) rows, a key, and constant bookkeeping
per event. Rebuilding filtered cache lists costs one such copy per
insertion scan. Even allocating a fresh binary cell for every arithmetic
temporary multiplies this quantity only by a polynomial in \(B\).
First count abstract records and scalar slots, treating each reference as one
slot independently of its eventual bit width. Give these slots consecutive
integer identifiers. Allocation indices therefore have length
\(O(\log S+\log(M+2)+\log(n+1)+\log(B+1))\).
But \(\log_2 S=\sum_{v\ne o}\log_2(H_v+1)=O(\ell)\),
and \(n\le\ell\), so this is \(O(B)\). Encoding each reference with that
width multiplies bit storage by only \(O(B)\), whose logarithm is still
\(O(B)\). This argument bounds the number
of allocations before charging their address lengths; it does not assume
the time bound it will be used to prove.
\end{proof}

\Needspace{22\baselineskip}
\begin{lemma}[Cost ledger]\label{lem:ledger}
The following conservative bounds include bit arithmetic and addressing:
\begin{center}
\small
\begin{tabular}{@{}p{.48\linewidth}p{.40\linewidth}@{}}
\toprule
Work & Bit-cost upper bound\\
\midrule
Parse and validate the explicit parent tree; normalize; compute \(k_0\);
clear rational denominators & \(O((n+1)^2B^3)\)\\
Enumerate all bounded vectors and test legality & \(O(S(n+1)^2B^3)\)\\
Lookup and insertion full-key comparisons, including list reconstruction
and possible copied answers & \(O((L+I)(M+1)(n+1)^2B^3)\)\\
Guards, residual subtraction, zero tests, calls, expansion and answer handling
& \(O((Q+C+E+1)(M+1)(n+1)^2B^3)\)\\
Serialize an explicit successful matrix & \(O((M+1)(n+1)^2B^3)\)\\
\bottomrule
\end{tabular}
\end{center}
\end{lemma}
\begin{proof}
Following parent pointers for at most \(n\) steps from every vertex checks
the rooted-tree condition in polynomial bit time, with rejection on an
invalid index or a cycle. Finite summation computes \(k_0\), and
Lemma~\ref{lem:bits} handles rational preprocessing.

The product of coordinate range sizes is exactly \(S\); fixing the root
to zero contributes one. In nested-loop enumeration, each prefix level
contains at most \(S\) nodes, since all remaining range sizes are at least
one. There are at most \((n+1)S\) nodes in total, including in the all-zero
case. Incrementing counters and materializing vectors costs polynomial
work in \(n,B\) per output. Each legality test checks bounds, a sum over
root children, and at most one parent condition per vertex. The scaled
cost uses at most \(n\) products and additions on \(O(B)\)-bit integers.
These give the stated (loose) enumeration and filtering bound.

Every full-key equality inspects the count and at most \(n\) residual
coordinates. It costs \(O((n+1)B^2)\) including accesses, with early
termination only reducing the cost. Both lookup and insertion filtering
perform this operation: their combined number is \(L+I\), not \(L\).
Insertion also traverses and rebuilds the filtered list. Charge each
visited/rebuilt entry to its insertion scan; charge the prepended entry
to the expansion. Even copying the stored answer per rebuilt entry is
covered by the \(M+1\) factor.

There is one coordinatewise guard per candidate check, at most one
materialized subtraction per accepted guard, and at most one all-zero
test per zero-depth expansion. Branch overhead, decrementing \(k\),
and empty-list tests are charged to calls and scans. A hit may copy an
answer of at most \(M\) rows, and a successful child may build one new
row list; allowing a full answer copy per such event is still within
the displayed bound. Lemma~\ref{lem:space} accounts for address lengths.
Finally output has at most \(M\) rows with \(n\) binary coordinates.
Neither physical-tour expansion nor optional certificate compression is
performed. All phases of the specified task have now been charged.
\end{proof}

\section{The resulting paper time theorem}\label{sec:time}
\begin{theorem}[Explicit-vector bit-cost implementation]\label{thm:time}
For the task in Definition~\ref{def:task}, there exists a deterministic
explicit-vector implementation of the ordered memoized search, with the
semantic correspondence of Lemma~\ref{lem:simulation}, whose time satisfies
\[
 \Time=O\bigl(S^3(M+1)(n+1)^2B^3\bigr)
      =O\bigl(S^3(M+1)(\ell+1)^5\bigr)
\]
in the specified bit-cost RAM model, for decision and explicit matrix output.
This is a paper theorem for that implementation, not a newly machine-checked
time theorem for the Lean executable.
\end{theorem}
\begin{proof}
Sum Lemma~\ref{lem:ledger} and apply \eqref{eq:logical}. Since \(S\ge1\),
\[
 1+S+C+Q+E+L+I=O(S^3).
\]
For example \(I\le S^2/2\), while the only cubic term is the lookup
bound. This gives the first estimate, including preprocessing and output.
Each input coordinate has at most \(\ell\) bits and \(n\le\ell\), so
\(M<n2^\ell\) when \(M>0\), and
\(\log_2(M+2)=O(\ell+1)\) also when \(M=0\).
Thus \(B=O(\ell+1)\), proving the second estimate.
\end{proof}

\begin{corollary}[Parameter classifications at the paper level]
The encoded decision and explicit matrix-witness tasks admit bit-cost RAM
upper bounds
\[
 O\bigl(2^{3M}(M+1)(\ell+1)^5\bigr)
       =2^{O(M)}(\ell+1)^{O(1)}
\]
and
\[
 O\bigl((U+1)^{3t}(tU+1)(\ell+1)^5\bigr).
\]
Consequently these tasks are FPT for \(M\) and for the joint parameter
\((t,U)\), for the paper implementation and representation specified here.
\end{corollary}
\begin{proof}
Use Lemma~\ref{lem:parameters} in Theorem~\ref{thm:time}.
The exponent of \(\ell+1\) is the fixed constant five. All other factors
depend only on the indicated parameter. In particular, \(M+1\le2^M\)
for natural \(M\), including zero. No XP intermediate result is needed.
\end{proof}

\begin{remark}[Other standard models]
These parameter classifications do not rely on free unbounded addressing.
For example, a deliberately naive multitape Turing-machine simulation has
polynomial overhead. If the charged RAM time is \(F\ge2\), at most \(F\)
memory operations occur and no written value or address has more than
\(O(F)\) bits. Store the touched cells as a sequential list of address-value
records, of total size \(O(F^2)\). Scanning it for each operation gives
\(O(F^3)\) time, with the already charged arithmetic performed on tapes.
This coarse simulation suffices for \(2^{O(M)}\ell^{O(1)}\) and FPT\((t,U)\)
on a standard tape model as well; it does not preserve the sharper cubic
power of \(S\) in Theorem~\ref{thm:time}.
\end{remark}

\begin{remark}[Limits of the conclusion]
No binary-input polynomial bound or XP bound for \(t\) alone follows.
On the one-edge tree, let \(H_v=U\), \(s=0\), and \(T=1\). Every legal
row has \(p_v=1\), so the explicit output has \(U=k_0\) rows although
\(t=1\) and the input length is \(O(\log(U+2))\). This output task already
requires \(\Omega(U)\) output symbols. It says nothing against a faster
decision algorithm or a directly constructed compressed certificate.
If the multiplicities are instead given in unary, \(M,U\le\ell\);
the displayed bound then gives an XP form in \(t\). None of these
statements yields a classification by tree height, treewidth, or vehicle
count alone, or a hardness reduction. Nor is \(2^{O(M)}\) a polynomial
in the numeric value \(M\).
\end{remark}

\section{What remains to formalize}
The paper proof fixes the computational task and discharges its mathematical
cost ledger. Following the verified component correspondences described above,
the next Lean package should prove exact ordered candidate-list correspondence,
then simulate the ordered search, attach a bit-cost semantics to its primitive
operations, and derive Theorem~\ref{thm:time}. The existing function-valued
search definitions have not been replaced in this package. The verified
contribution includes logical counts, parameters, representation components,
and successful-cache-entry bounds; the additional time and FPT statements
above still have paper-proof status only.

\end{document}
